\newpage
\section{Diskussion}

\subsection{Interpretation der Ergebnisse}

Der steigende \ac{UI}-Kontext ist nicht allein für die höhere Testqualität verantwortlich. Der Accessibility-Snapshot aus Stufe~2 hilft bei der Behebung der Selektorenprobleme. Erfundene \code|data-testid|-Werte fallen von 39,4\,\% auf 0,4\,\% der Dateien, der \emph{selector}-Median steigt von 2 auf 3 und der \texttt{INFRA\_FAIL}-Anteil sinkt von 69,0\,\% auf 42,2\,\%. Die Map-Model-Helfer behandeln ein anderes Problem, nämlich den Zugang zu Zuständen, die nicht im \ac{DOM} liegen. Die Map-Model-Helfer in Stufe~3 sorgen dafür, dass \emph{map\_interaction} im Mittel um 1,09 Punkte steigt und die PASS-Rate um 13,7 Prozentpunkte ansteigt. Der Anstieg der PASS-Rate ist aber auch auf den erstmaligen vollständigen \ac{UI}-Kontext der Webanwendung zurückzuführen. Beide Sprünge treten in anderen Stufen und anderen Dimensionen auf. Dem Modell fehlt demnach Wissen darüber, welche Elemente existieren und wie sie benannt sind, und ein Weg, den Kartenzustand überhaupt abzufragen. In Stufe~1 ist zudem auch sichtbar, dass die Tests nicht scheitern, weil das \ac{LLM} Playwright nicht beherrscht, sondern weil es die Anwendung nicht kennt.

Der zweite Befund betrifft die Aussagekraft des Ausführungsergebnisses. In Stufe~1 sind 33\,\% der bestandenen Tests als \texttt{vacuous\_pass} eingestuft, nahezu alle davon bei UC-03. Ohne die Map-Model-Helfer werden nur Sichtbarkeiten der Elemente im \ac{DOM} geprüft, nicht aber das im Use Case geforderte Zoomverhalten, genau deshalb ist die PASS-Rate auch so hoch. Ab Stufe~3 versuchen die Tests die richtige Prüfung und scheitern an der Umsetzung, wodurch die PASS-Rate deutlich sinkt. Der Rückgang ist allerdings kein Qualitätsverlust, sondern das Gegenteil. Eine reine Bewertung nach der PASS-Rate hätte für Stufe~1 das beste Ergebnis ausgewiesen, obwohl die Prüfung nicht die ist, nach der verlangt wurde. Die Evaluation aus Phase~2 zeigt also, dass das eigentliche Kriterium nicht ist, ob ein Test besteht, sondern ob er bei einem tatsächlichen Fehler in der Anwendung auch wirklich fehlschlagen würde.

Die manuell erstellte \ac{UI}-Map in Stufe~4 enthält das vollständige Wissen der Struktur, macht aber keinen großen Unterschied zu Stufe~3. Die häufigste Fehlerkategorie bleibt \texttt{ASSERTION\_FAIL} mit 30,8\,\%. Welches Element aufzufinden ist, lässt sich auch aus einer einfachen Elementliste ableiten, aber wie der erwartete Zustand genau zu prüfen ist, hingegen nicht. Die Übergangsraten aus Stufe~5 verdeutlichen dies: Die Übergangsrate von einem \texttt{INFRA\_FAIL} zu einem bestandenen Test ist höher als aus einem \texttt{ASSERTION\_FAIL} und 83,3\,\% der Fehler aus \texttt{ASSERTION\_FAIL} bleiben auch in dieser Kategorie. Der Fehler bei der Weiterverwendung von \code|expect.poll| wird auch durch den Loop nicht vollständig behoben, da er nichts mit dem Kontext zu tun hat. Stattdessen sollte die Information im OPT-Playwright-Skill ergänzt werden.

Der bereitgestellte \ac{UI}-Kontext ist in Stufe~4 noch nicht an seinem Maximum angekommen. Was noch fehlt, ist Wissen über Lage, Größe und Überlagerung der Elemente. In Stufe~3 klicken 24 der 50 UC-06-Dateien auf die Position (100, 100) und treffen dort den Layer Switcher statt der Karte. Bei UC-08 verdeckt das Messpanel den Bereich, in dem die Messlinie gezeichnet wird. Bei UC-07 ist die Umrechnung von Karten- in Pixelkoordinaten über eine passende Funktion überhaupt nicht lösbar. 

Die Bonus-Stufe~5 hat mit 73,2\,\% die höchste PASS-Rate, benötigt dafür aber 2.154 statt 500 Generierungsaufrufe und etwa die vierfache Laufzeit. Der Nutzen liegt aber deutlich in der Korrektur der ersten Iterationen. Einige Tests scheitern nach der ersten Generierung an Kleinigkeiten, die in den ersten Iterationen durch den neuen Kontext und die Fehlermeldung von Playwright direkt behoben werden können. Falsch gewählte Prüfstrategien und Karteninteraktionen können dadurch aber nicht abgeleitet und verbessert werden. Der zusätzliche Screenshot dient zwar als visuelle Hilfe, reicht aber nicht immer aus. Außerdem ist der Loop so konfiguriert, dass er bei einem \texttt{PASS} abbricht und damit auf ein Kriterium optimiert, das mit der Testqualität nur bedingt zusammenhängt. Die Gefahr für einen \texttt{vacuous\_pass} steigt damit. In dieser Arbeit ist das allerdings kein Problem, da die manuellen Referenztests belegen, dass die Anwendung funktioniert und ein erfolgreicher Test daher plausibel ist. In einem produktiven Umfeld, in dem getestet wird, ob tatsächlich Fehler in der Anwendung sind, kann das zu einem Problem werden.

Die Bewertung der Tests durch LLM-as-a-Judge bedarf ebenfalls einer kurzen Einordnung. Der in Abschnitt~6.3.1 beschriebene Regelkatalog war im Prompt nicht vorgegeben, wurde aber in getrennten Kontextfenstern wiederholt verwendet und nicht durchgängig gleich angewendet. Die Punkte aus der Bewertung sind daher nicht zu 100\,\% mit den bereitgestellten Rubriken in Anhang~\ref{app:judge-rubrik} reproduzierbar. Der Grund liegt aber wohl eher in der Rubrik selbst, da nicht für alle Fälle klar definiert war, wie das \ac{LLM} die Bewertung einordnen soll. Die Bestätigungsquote der Stichprobenkontrolle steigt von 80,0\,\% in Stufe~1 auf 97,1\,\% in Stufe~5, demnach erzeugen bessere Tests weniger Grenzfälle.

Der zweite Generierungslauf mit GPT-5.4 zeigt, dass die Auswirkungen des \ac{UI}-Kontexts nicht an das verwendete Modell gebunden sind. Die Übergänge zwischen den Stufen treten dort ebenfalls auf. Auch hier verschwinden erfundene Test-IDs von Stufe~1 auf 2, von Stufe~2 auf 3 steigt \emph{map\_interaction} im Mittel von 2,00 auf 3,75 und die PASS-Raten steigen über die Stufen 1 bis 3 an. Die Testqualität ist also abhängig von der Form des bereitgestellten \ac{UI}-Kontexts. Stufe~4 fällt allerdings gegenüber Stufe~3 etwas schlechter aus, was noch einmal verdeutlicht, dass der Unterschied zwischen diesen beiden Stufen nicht sonderlich groß ist.

\subsection{Limitationen}

\subsection{Handlungsempfehlungen}
